\documentclass{article}
\usepackage{graphicx} % Required for inserting images
\usepackage{geometry}
\usepackage{amsmath}
 \geometry{
 a4paper,
 total={190mm,257mm},
 left=10mm,
 top=20mm,
 }

\title{Quantum Management\\An Analogy between Management and Physics}
\author{}
\date{}

\begin{document}
\begin{description}
    \item[Candidate: Francesca Lonigro]
    \item[Company Tutors: Dr. Anna Deambrosis \& Dr. Davide Bugini]
    \item[Academic Supervisor: Prof. Elisa Luciano]
\end{description}
\subsubsection*{\centering{Quantum Management - An Analogy between Management and Physics}}
\begin{abstract}
    This thesis explores a novel management theory, relativistic quantum management, building on Danah Zohar's concept of quantum management. While Zohar's theory embraces uncertainty and indeterminism, this work extends it by introducing the principle of relativity, drawing an analogy between management and quantum field theory. I propose a mathematical formalization of relativistic quantum management using frameworks from physics: Feynman path integrals to model the phase dependence of a process, De Broglie-Bohm Pilot wave theory to represent the project backlog, the Von Neumann Density matrices to describe the sprint backlog and Einstein Special relativity to account for the perception of spacetime within interaction frameworks. To simulate this theoretical model, I developed a neural network called $\phi \epsilon \rho \omega$ Quantum Intelligence Neural Network ($\phi QINN$). I applied this model to various insurance domains. The thesis highlights also how relativistic quantum management provides a structured approach to decision-making under uncertainty, referencing the theories of Kahneman and Tversky. In conclusion, the $\phi QINN$ simulation shows a management approach that navigates information theory (quantum or classical) and thrives on the inherent error, risk, uncertainty, indeterminism and relativity present in a process of a complex system.
\end{abstract}
\subsubsection*{Takeaways}
\begin{enumerate}
    \item \textbf{$\phi \epsilon \rho \omega$} and its \textbf{Neural Network}:
    Relativistic Quantum Management is thought as an extension of Quantum Management, by including Special Relativity.
    The central topic of the thesis is the likely propagator
     \begin{equation*}
    \phi \epsilon \rho \omega = \int \mathcal{D}scenario \exp\left( \frac{i}{\text{\textbf{Quantum Term}}} Action[scenario] \right)
    \end{equation*}
    %This is in particular inspired by the Feynman Path Integral and it is envisioned as multiple paths weighted by a phase and a functional integral that sums all up, a pilot wave that is the one that says from where to start and where to end as a project backlog, Von Neumann density matrices represent the intermediate states or sprint backlog of a project. Thus, from the beginning to the end, there could be a multitude of paths. The way I decided to weight the best one is by considering the entropy. The one with the lowest entropy is the best trajectory, and thus, the best trajectory of decisions that someone could do in a given step of a process.\end{comment}\\
    The \textbf{quantum term} has been thought so that if it goes to zero, as a physical classical limit, it brings the quantum management back to the Newtonian one.
    \item \textbf{The Quantum Term \textit{Adaptiveness} in Leadership}: Adaptiveness is a quantum factor that lets an organization sense, respond, and learn to find the most efficient path for a project. Measuring adaptiveness requires new metrics focused on processes like the speed of response and learning rate, rather than just the final output. In this model, a project's success is tied to how well the organization navigated all possible scenarios, acting as a dynamic and adaptive system
    \item \textbf{The Quantum Term \textit{Realism} in Asset Management}: Realism is a leader's ability to base decisions on market complexities using a probabilistic framework. Measuring realism involves adopting a non-deterministic and  deterministic model, such as $\phi \epsilon \rho \omega$, to account for a wide range of market states, including low-probability "tail risks" and "Black Swan" events. Metrics for realism include the number of scenarios considered and how well a model's simulated return distribution matches the skewness and kurtosis of historical data.
    \item \textbf{The Quantum Term \textit{Balance} in Environmental, Social and Governance Principles}: Balance in ESG means harmoniously considering environmental, social, and governance factors rather than prioritizing one over the others. Measuring ESG principles and their balance requires a mix of quantitative and qualitative metrics that go beyond just financial performance. Using quantum computing can help leaders model this complex system to proactively maintain balance and make better decisions.
    \item \textbf{The Quantum Term \textit{Proactivity} in Customer Centricity}: Proactivity, a "quantum" term in the Path Integral model, is an organization's ability to anticipate future needs, to act efficiently on existing ones by also learning from the past. It can be measured by an Early Warning Indicators (EWI) score, which tracks how often a company successfully identifies emerging trends before they become widely known. Other metrics include the accuracy of customer needs predictions and the average time for predictive product development
    \item \textbf{Information Theory in Relativistic Quantum Management}: Information theory is the mathematical study of the quantification, storage, and communication of information. This can be classical or quantum and it can be in relativistic or non-relativistic mode.

    The psychological information processing can summarized in four main steps:
    \textbf{encoding, storage, retrieval and transformation}. At the core of transformation there is the decision-making process. 
    \begin{equation*}
        \text{How do we make a decision?}
    \end{equation*}
    The Relativistic Quantum Management embraces the uncertainty, the indeterminism and the relativity in this process by also considering the existence of bias and noise. It highlights how the information processing is thus influenced by the same bias and noise.

\end{enumerate}
\end{document}
